\section*{ВВЕДЕНИЕ}
\addcontentsline{toc}{section}{ВВЕДЕНИЕ}
\thispagestyle{plain}

\textbf{Актуальность работы.}
Цифровая трансформация высшего образования и распространение
гибридных и дистанционных форматов обучения формируют запрос
на инструменты, обеспечивающие активное вовлечение студентов,
мгновенную обратную связь преподавателю и измеримое повышение
качества усвоения учебного материала~\cite{sagindykova2015,braude2020}.
Классические лекции и семинары всё чаще дополняются интерактивными
сессиями: опросами в реальном времени, викторинами с подсчётом очков,
сессиями вопросов--ответов с возможностью голосования за наиболее
важные вопросы аудитории.

Анализ существующих решений (Quizlet, Kahoot, Mentimeter, Slido)
показывает, что при всех функциональных возможностях они опираются
на инфраструктуру иностранных облачных провайдеров.
В условиях функционирования российских образовательных учреждений
такая зависимость снижает надёжность сервиса, ограничивает контроль
над персональными данными обучающихся и в ряде случаев делает
недоступной полную функциональность.
В отечественных аналогах (JoyTeka, OnlineTestPad) геймификация,
как правило, ограничена баллами и таблицей лидеров,
а возможности интеграции с университетскими LMS реализованы
фрагментарно.

В предшествующей научно-исследовательской работе автором был
разработан прототип сервиса EduQuiz на стеке Flutter и Firebase.
Прототип позволил отработать сценарии интерактивного взаимодействия,
однако выявил существенные ограничения: зависимость от инфраструктуры
Google, минимальную глубину геймификации, отсутствие интеграции
с университетской системой управления обучением (Learning
Management System, LMS) \textit{edu.gubkin}, а также отсутствие
расчётного обоснования показателей надёжности и экспериментальной
проверки педагогической эффективности.
Перечисленные ограничения определяют необходимость переработки
архитектуры с переносом backend-логики на собственный
серверный компонент, развёртываемый в инфраструктуре Российской
Федерации, и сопровождением полученного решения расчётами надёжности
и экспериментальным исследованием. Указанные обстоятельства
обусловливают актуальность настоящей магистерской работы.

\textbf{Объект исследования} --- процессы интерактивного
взаимодействия преподавателя и студентов в ходе аудиторных
и дистанционных занятий в высшем учебном заведении.

\textbf{Предмет исследования} --- модели, алгоритмы и программные
средства организации интерактивного тестирования с элементами
геймификации, обеспечивающие требуемые показатели надёжности
сервиса и измеримое повышение качества усвоения учебного материала.

\textbf{Цель работы} --- разработка моделей и сервиса
для интерактивного взаимодействия с обучающимися, обеспечивающего
проведение интерактивных сессий с расширенной геймификацией,
импорт банка вопросов из университетской LMS \textit{edu.gubkin}
и работу в инфраструктуре Российской Федерации с расчётно
обоснованными показателями надёжности.

\textbf{Задачи работы.} Для достижения поставленной цели
в работе решаются следующие задачи:
\begin{enumerate}
    \item Провести анализ существующих систем интерактивного
          обучения, теоретических подходов к геймификации
          и интервальному повторению, а также положений
          теории надёжности применительно к разрабатываемому сервису.
    \item Разработать модели взаимодействия с обучающимися:
          ролевую модель пользователей, модели курса и банка
          вопросов, модель интерактивной сессии, модель игровой
          прогрессии, модель интервального повторения по алгоритму
          SM\textendash 2 и модель оценки надёжности усвоения знаний.
    \item Спроектировать клиент-серверную архитектуру сервиса
          с заменой стека Firebase на собственный backend
          на языке Go, реализовать кроссплатформенный клиент
          на Flutter (Web, iOS, Android, macOS), реализовать
          интеграцию с университетской LMS \textit{edu.gubkin}.
    \item Выполнить расчёт показателей надёжности сервиса
          (вероятности безотказной работы, средней наработки на отказ,
          коэффициента готовности), обосновать стратегию резервирования
          и провести нагрузочное тестирование; реализовать меры
          информационной безопасности.
    \item Разработать методику педагогического эксперимента
          со студентами курса \textit{«Надёжность»} и реализовать
          программный пайплайн обработки результатов с применением
          методов математической статистики.
\end{enumerate}

\textbf{Научная новизна} результатов работы заключается в:
\begin{enumerate}
    \item переходе от рассмотренной в предшествующей научно-исследовательской
          работе serverless-инфраструктуры Firebase к собственному
          backend-компоненту, развёртываемому в инфраструктуре
          Российской Федерации, с расчётным обоснованием показателей
          надёжности предложенной архитектуры;
    \item объединении интерактивной учебной сессии и индивидуального
          интервального повторения по алгоритму SM\textendash 2
          через автоматическое формирование оценки качества ответа
          по паре «корректность ответа --- время ответа», без явного
          запроса оценки у обучающегося;
    \item связи учебного содержания и предмета разработки в рамках
          одной работы: положения теории надёжности, преподаваемой
          на кафедре автоматизированных систем управления РГУ Нефти
          и Газа (НИУ) им. И.М. Губкина, применены к расчёту показателей
          надёжности собственной архитектуры сервиса, выступающего
          средством преподавания этой же дисциплины.
\end{enumerate}

\textbf{Методы исследования.}
В работе использованы методы системного анализа,
объектно-ориентированного проектирования (UML, нотация C4),
теории надёжности, классической теории тестов (коэффициент
$\alpha$ Кронбаха, индекс трудности пунктов, точечно-бисериальная
корреляция), а также методы математической статистики (критерии
Шапиро\textendash Уилка, Стьюдента, Манна\textendash Уитни,
оценка размера эффекта по $d$ Коэна, поправка Бонферрони).

\textbf{Практическая значимость.}
Разработанный сервис развёртывается в инфраструктуре Российской
Федерации и выпускается в виде кроссплатформенного клиента
(Web, iOS, Android, macOS), что делает его применимым к преподаванию
любых дисциплин в РГУ Нефти и Газа им. И.М. Губкина. В качестве
содержательной площадки выбран курс \textit{«Надёжность»}, для
которого подготовлен демонстрационный банк вопросов и разработана
методика педагогического эксперимента; апробация на реальной группе
студентов запланирована в текущем семестре.

\textbf{Апробация результатов.}
Апробация инструментальной части работы выполнена в форме нагрузочного
тестирования собственным программным инструментом при $N = 30, 100, 300$
и $500$ одновременных подключениях (см. главу~3); функциональная
апробация проведена на демонстрационном банке вопросов по теории
надёжности. Апробация педагогической методики на студентах курса
\textit{«Надёжность»} кафедры автоматизированных систем управления
факультета автоматики и вычислительной техники РГУ Нефти и Газа
(НИУ) им. И.М. Губкина запланирована в весеннем семестре 2025/2026
учебного года.

\textbf{Положения, выносимые на защиту.}
\begin{enumerate}
    \item Архитектура сервиса интерактивного тестирования с собственным
          backend на языке Go и кроссплатформенным клиентом на Flutter,
          развёртываемая в инфраструктуре Российской Федерации.
    \item Модели игровой прогрессии и интервального повторения,
          объединённые через автоматическое формирование оценки качества
          ответа по паре «корректность ответа --- время ответа»
          и применённые к содержанию дисциплины \textit{«Надёжность»}.
    \item Расчёт показателей надёжности разработанной архитектуры
          (вероятности безотказной работы, средней наработки на отказ,
          коэффициента готовности) для конфигурации без резервирования
          и с горячим резервированием СУБД.
    \item Методика педагогического эксперимента и реализованный
          программный пайплайн её обработки методами математической
          статистики; готовность к апробации на реальной группе
          студентов курса \textit{«Надёжность»}.
\end{enumerate}

\textbf{Структура работы.}
Магистерская диссертация состоит из введения, трёх глав, заключения,
списка использованных источников и приложений; объём основного
текста 60--70 страниц, общий объём с приложениями --- до 100 страниц.
В первой главе проводится анализ предметной области и формулируется
постановка задачи. Во второй главе разрабатываются модели взаимодействия
с обучающимися и проектируется архитектура сервиса. В третьей главе
обосновывается надёжность и безопасность архитектуры и описывается
экспериментальное исследование. В заключении формулируются основные
итоги работы.

\newpage
